\section{A Joint Occurrence--Detection Model}
\subsection{Observations and Assumptions}
For account $i$, let $T_{f,i}$ be its first fraud time and $T_{d,i}$ its first
platform detection time. We observe
\[
 y_i=\min(T_{d,i},C_i),\qquad
 \delta_i=\mathbf{1}\{T_{d,i}\le C_i\},
\]
where $C_i$ is an administrative observation boundary. We also observe a
timestamped covariate history $\Hist_i(t)$ up to $y_i$. Time zero is the start
of the modeled risk period. Initially, all accounts are assumed fraud-free;
prevalent fraud at enrollment requires an additional initial-state mixture.

The initial model assumes that detection follows genuine fraud, with
$T_d\ge T_f$, and that fraud status is absorbing. False suspensions,
repeated fraud episodes, and rehabilitation are outside this first model.
Censoring is assumed non-informative conditional on the modeled information.
These are modeling assumptions to audit, not established properties of the
Twitter or Wiki data.

We first specify the likelihood conditional on an external covariate path;
path conditioning is suppressed in the notation. Neural hazards may use
causal features available immediately before each transition. If covariates
are themselves generated by the latent fraud state, this conditional
construction needs further justification or a joint observation model.
No future observations may be used for an online prediction.

\subsection{Continuous-Time Likelihood}
Let $\lambda_f(t)$ be the hazard of first fraud, and let
$\lambda_d(t\mid u)$ be the detection hazard at $t$ after fraud at $u$.
Dependence on onset $u$ can include elapsed time $t-u$. Define
\begin{align}
 S_f(t)&=\exp\left[-\int_0^t\lambda_f(s)\,ds\right],\\
 f_f(u)&=\lambda_f(u)S_f(u),\\
 S_d(t\mid u)&=\exp\left[-\int_u^t\lambda_d(s\mid u)\,ds\right],\\
 f_d(t\mid u)&=\lambda_d(t\mid u)S_d(t\mid u).
\end{align}
We do not require that all accounts eventually commit fraud. The
occurrence distribution may retain probability mass at infinity.

For a detection observed at $d$, the marginal density is
\begin{equation}
 p_d(d)=\int_0^d f_f(u)f_d(d\mid u)\,du.
 \label{eq:detected}
\end{equation}
For an account undetected at $C$, its probability is
\begin{equation}
 P(T_d>C)=S_f(C)+\int_0^C f_f(u)S_d(C\mid u)\,du.
 \label{eq:censored}
\end{equation}
The terms in Equation~\ref{eq:censored} represent no fraud by $C$, and fraud
by $C$ without detection, respectively. Treating all undetected users as
fraud-free removes the second explanation.
The observed-data negative log-likelihood is
\begin{equation}
 \Loss=-\sum_i\left[\delta_i\log p_d(y_i)
 +(1-\delta_i)\log P(T_{d,i}>y_i)\right].
 \label{eq:nll}
\end{equation}
Exact timestamps contribute densities. If platform status is checked only
periodically, detection is interval-observed and its contribution is instead
$P(l_i<T_{d,i}\le r_i)$. This distinction must be resolved in preprocessing.

\subsection{Discrete-Time Specification}
For an initial implementation, let $k=1,\ldots,K$ index observation bins.
Let $a_k\in(0,1)$ be the probability of first fraud in bin $k$, conditional on
no previous fraud. Given onset in bin $k$, let $b_{j\mid k}\in(0,1)$ be the
conditional detection probability in bin $j\ge k$, given no earlier detection.
The bin convention allows onset followed by detection within the same bin.
This defines a discrete model; it is not an exact continuous-time model
without an additional discretization argument.

The onset mass and conditional detection survival are
\begin{align}
 g_k&=a_k\prod_{r<k}(1-a_r),\\
 Q_{j\mid k}&=\prod_{r=k}^{j}(1-b_{r\mid k}).
\end{align}
An empty product is one. Detection in bin $d$ and non-detection through bin
$c$ have likelihoods
\begin{align}
 p_d(d)&=\sum_{k=1}^{d}g_k Q_{d-1\mid k}b_{d\mid k},
 \label{eq:discrete-detected}\\
 p_{\mathrm{und}}(c)&=\prod_{r=1}^{c}(1-a_r)
                  +\sum_{k=1}^{c}g_k Q_{c\mid k}.
 \label{eq:discrete-censored}
\end{align}
These finite sums permit exact marginalization over onset bins and direct
automatic differentiation, avoiding a sampling approximation in the first
experiment. Computation should use log probabilities and log-sum-exp.
General onset-dependent detection costs $O(K^2)$ per sequence.

\paragraph{A linear-time restricted model.}
If detection depends on current causal features but not onset bin,
$b_{j\mid k}=b_j$, define $U_k$ as the probability of no fraud through bin $k$
and $V_k$ as the probability of fraud but no detection through that bin.
Starting with $U_0=1$ and $V_0=0$,
\begin{align}
 U_k&=U_{k-1}(1-a_k),\\
 V_k&=(V_{k-1}+U_{k-1}a_k)(1-b_k),\\
 D_k&=(V_{k-1}+U_{k-1}a_k)b_k.
\end{align}
Here $D_k$ is the probability mass of first detection in bin $k$.
The censored likelihood is $U_c+V_c$. These recurrences imply
$U_c+V_c+\sum_{k=1}^cD_k=1$, providing a useful implementation check.
A shared causal GRU with separate sigmoid heads is a candidate
parameterization of $a_k$ and $b_k$, not a commitment that this
parameterization resolves identification.

\subsection{Retrospective and Online Inference}
For a detected account, retrospective onset responsibilities are
\begin{equation}
 P(T_f=k\mid T_d=d)
 =\frac{g_k Q_{d-1\mid k}b_{d\mid k}}{p_d(d)}.
\end{equation}
For an undetected account, the posterior includes both an onset-bin
component $g_kQ_{c\mid k}/p_{\mathrm{und}}(c)$ and a no-fraud-yet component
$U_c/p_{\mathrm{und}}(c)$.
In the restricted model, an online score among undetected accounts is
\begin{equation}
 R_c=P(T_f\le c\mid T_d>c,\Hist(c))=\frac{V_c}{U_c+V_c}.
\end{equation}
This filtered probability need not be monotone: continued non-detection can
change the evidence about past fraud. An absorbing alert policy may retain a
previous alert, but that policy is distinct from probability monotonicity.
Retrospective smoothing uses later detection information and must not be
reported as an online prediction made earlier.
